\SourcePage{781}{784}
\renewcommand{\thesection}{\alph{section}}
\setcounter{section}{2}
\section[Legendre polynomials]{Legendre polynomials\protect\footnote{The mathematical literature contains many good accounts of the theory of spherical functions. Here, for reference, we give only some basic relations, without attempting any systematic presentation of the theory of these functions.}}

The \emph{Legendre polynomials} $P_l(\cos\theta)$ are defined by
\begin{equation}
 P_l(\cos\theta)=\frac1{2^ll!}
 \frac{d^l}{(d\cos\theta)^l}(\cos^2\theta-1)^l.
 \tag{c.1}
\end{equation}
They satisfy the differential equation
\begin{equation}
 \frac1{\sin\theta}\frac d{d\theta}
 \left(\sin\theta\frac{dP_l}{d\theta}\right)+l(l+1)P_l=0.
 \tag{c.2}
\end{equation}

\SourcePage{782}{785}
The associated Legendre polynomials are defined by
\begin{equation}
 \begin{split}
 P_l^m(\cos\theta)
 &=\sin^m\theta\frac{d^mP_l(\cos\theta)}{(d\cos\theta)^m}\\
 &=\frac1{2^ll!}\sin^m\theta
 \frac{d^{l+m}}{(d\cos\theta)^{l+m}}(\cos^2\theta-1)^l
 \end{split}
 \tag{c.3}
\end{equation}
or, equivalently, by
\begin{equation}
 P_l^m(\cos\theta)=(-1)^m\frac{(l+m)!}{(l-m)!2^ll!}
 \sin^{-m}\theta\frac{d^{l-m}}{(d\cos\theta)^{l-m}}
 (\cos^2\theta-1)^l,
 \tag{c.4}
\end{equation}
where $m=0,1,\ldots,l$. The associated polynomials obey the equation
\begin{equation}
 \frac1{\sin\theta}\frac d{d\theta}
 \left(\sin\theta\frac{dP_l^m}{d\theta}\right)
 +\left[l(l+1)-\frac{m^2}{\sin^2\theta}\right]P_l^m=0.
 \tag{c.5}
\end{equation}

The normalization integral of the Legendre polynomials,
$\int_{-1}^{1}[P_l(\mu)]^2d\mu$ ($\mu=\cos\theta$), is evaluated by
substituting (c.1) into it and integrating by parts $l$ times. It then becomes
\[
 \frac{(-1)^l}{2^{2l}(l!)^2}\int_{-1}^{1}(\mu^2-1)^l
 \frac{d^{2l}}{d\mu^{2l}}(\mu^2-1)^l\,d\mu
 =\frac{(2l)!}{2^{2l}(l!)^2}\int_{-1}^{1}(1-\mu^2)^l\,d\mu.
\]
With the substitution $u=(1-\mu)/2$, this integral reduces to Euler's
$B$ integral and gives
\begin{equation}
 \int_{-1}^{1}[P_l(\mu)]^2d\mu=\frac2{2l+1}.
 \tag{c.6}
\end{equation}
In the same way, one readily verifies that the functions $P_l(\mu)$ with
different $l$ are mutually orthogonal:
\begin{equation}
 \int_{-1}^{1}P_l(\mu)P_{l'}(\mu)\,d\mu=0,\hspace{1.5em}l\ne l'.
 \tag{c.7}
\end{equation}

The normalization integral for the associated polynomials can be calculated
similarly: write $[P_l^m(\mu)]^2$ as the product of expressions (c.3) and
(c.4), and integrate by parts $l-m$ times. The result is
\begin{equation}
 \int_{-1}^{1}[P_l^m(\mu)]^2d\mu
 =\frac2{2l+1}\frac{(l+m)!}{(l-m)!}.
 \tag{c.8}
\end{equation}

\SourcePage{783}{786}
It is likewise easy to verify that functions $P_l^m$ with different $l$ (and
the same $m$) are mutually orthogonal:
\begin{equation}
 \int_{-1}^{1}P_l^m(\mu)P_{l'}^m(\mu)\,d\mu=0,
 \hspace{1.5em}l\ne l'.
 \tag{c.9}
\end{equation}

Integrals of products of three Legendre polynomials were considered in
\S~107.

The following \emph{addition theorem} holds for the Legendre polynomials. Let
$\gamma$ be the angle between two directions specified by the spherical
angles $\theta,\varphi$ and $\theta',\varphi'$:
$\cos\gamma=\cos\theta\cos\theta'+\sin\theta\sin\theta'
\cos(\varphi-\varphi')$. Then
\begin{equation}
 \begin{split}
 P_l(\cos\gamma)&=P_l(\cos\theta)P_l(\cos\theta')\\
 &\hspace{1.5em}+\sum_{m=1}^{l}2\frac{(l-m)!}{(l+m)!}
 P_l^m(\cos\theta)P_l^m(\cos\theta')
 \cos[m(\varphi-\varphi')].
 \end{split}
 \tag{c.10}
\end{equation}
In terms of the spherical functions (defined as in (28.7)), this theorem may
also be written
\begin{equation}
 P_l(\mathbf n\mathbf n')=\frac{4\pi}{2l+1}
 \sum_{m=-l}^{l}Y_{lm}^*(\mathbf n')Y_{lm}(\mathbf n).
 \tag{c.11}
\end{equation}
Here $\mathbf n$ and $\mathbf n'$ are two unit vectors, and
$Y_{lm}(\mathbf n)$ denotes the spherical function of the polar angle and
azimuth of the direction $\mathbf n$ relative to a fixed coordinate system.

Multiply (c.10) by $P_{l'}(\cos\theta)$ and integrate it over
$do=\sin\theta\,d\theta\,d\varphi$. Integration with respect to $d\varphi$
makes all terms on the right-hand side containing factors
$\cos[m(\varphi-\varphi')]$ vanish. Taking account of (c.6) and (c.7), we
obtain
\[
 \int P_l(\cos\gamma)P_{l'}(\cos\theta)\,do
 =\delta_{ll'}\frac{4\pi}{2l+1}P_l(\cos\theta').
\]
This result can be written in the symmetric form
\begin{equation}
 \int P_l(\mathbf n_1\mathbf n_2)P_{l'}(\mathbf n_1\mathbf n_3)\,do_1
 =\delta_{ll'}\frac{4\pi}{2l+1}P_l(\mathbf n_2\mathbf n_3),
 \tag{c.12}
\end{equation}
where $\mathbf n_1$, $\mathbf n_2$, and $\mathbf n_3$ are three unit vectors,
and the integration is performed over the directions of one of them,
$\mathbf n_1$. Finally,

\SourcePage{784}{787}
we list the first few normalized spherical functions $Y_{lm}$:
\begin{gather*}
 Y_{00}=\frac1{\sqrt{4\pi}},\hspace{2em}
 Y_{10}=i\sqrt{\frac3{4\pi}}\cos\theta,\\
 Y_{1,\pm1}=\pm i\sqrt{\frac3{8\pi}}\sin\theta\,e^{\pm i\varphi},\\
 Y_{20}=\sqrt{\frac5{16\pi}}(1-3\cos^2\theta),\\
 Y_{2,\pm1}=\pm\sqrt{\frac{15}{8\pi}}
 \cos\theta\sin\theta\,e^{\pm i\varphi},\\
 Y_{2,\pm2}=-\sqrt{\frac{15}{32\pi}}\sin^2\theta\,e^{\pm2i\varphi},\\
 Y_{30}=-i\sqrt{\frac7{16\pi}}\cos\theta(5\cos^2\theta-3),\\
 Y_{3,\pm1}=\pm i\sqrt{\frac{21}{64\pi}}
 \sin\theta(5\cos^2\theta-1)e^{\pm i\varphi},\\
 Y_{3,\pm2}=-i\sqrt{\frac{105}{32\pi}}
 \cos\theta\sin^2\theta\,e^{\pm2i\varphi},\\
 Y_{3,\pm3}=\pm i\sqrt{\frac{35}{64\pi}}
 \sin^3\theta\,e^{\pm3i\varphi}.
\end{gather*}
